% ---------------------------------------------------------------------------
% Conclusiones y recomendaciones (norma 3.14)
% ---------------------------------------------------------------------------
% NO se numera como capítulo (norma 2.6) y sus elementos NO deben numerarse.
%
% La norma pide:
%   - Conclusiones derivadas de forma lógica del trabajo y sustentadas por los
%     resultados obtenidos.
%   - Recomendaciones bien razonadas y fundamentadas.
%   - Pueden incluirse sugerencias sobre la aplicación de los resultados y
%     perspectivas para trabajos futuros.

\seccionSinNumerar{Conclusiones y recomendaciones}

El objetivo general de este trabajo era desarrollar un prototipo que clasificara y priorizara
alertas de seguridad, justificando cada decisión con un modelo de lenguaje apoyado en generación
aumentada por recuperación ejecutada en local, y que redujera los falsos positivos sin sacrificar
la detección de amenazas reales. Los resultados del capítulo \ref{cap:resultados} sostienen que ese
objetivo se cumplió dentro del caso de uso acotado sobre el que se evaluó: el prototipo suprimió
los falsos positivos que el nivel de regla de Wazuh deja pasar ---tasa de 0,000 frente a 0,282---,
sin dejar de detectar ninguna amenaza real y sin ejecutar ninguna acción que hubiera interrumpido
indebidamente un servicio del cliente. Este resultado confirma, con datos y no solo con diseño, la
idea central del trabajo: separar la clasificación, la decisión y la justificación en componentes
distintos ---una política determinista que decide, un perfil de cliente que filtra, y un modelo de
lenguaje que solo explica--- permite obtener las ventajas de un modelo de lenguaje sin heredar su
principal debilidad, que es decidir de forma no auditable.

La conclusión metodológica más relevante del trabajo se refiere a la medida de la explicabilidad:
\textbf{un anclaje verificable perfecto no implica una justificación correcta}, y ambas propiedades
deben medirse por separado. La verificación de anclaje comprueba que el modelo no inventa campos, y
eso es necesario, pero se puede satisfacer sin explicar nada. Se documentaron tres maneras
distintas de superarla vacíamente: describir al revés el significado de los campos citados, repetir
el enunciado recibido ---que contiene por construcción todos los campos--- y, la más
problemática, contradecir la clasificación que el propio sistema acaba de emitir. Una métrica
automática de explicabilidad debe defenderse de las formas de cumplirla sin cumplir su propósito.

De ahí se sigue una lectura más matizada de la recuperación aumentada que la que el trabajo
esperaba encontrar. Medida por clase, la recuperación mejoraba con claridad la justificación de las
amenazas reales y empeoraba la de los falsos positivos, hasta el punto de contradecir la
clasificación que el propio motor acababa de emitir. La causa no estaba en la recuperación sino en
el enunciado, que preguntaba «por qué importa esta alerta» incluso cuando el sistema acababa de
decidir que no importa; corregido eso, la contradicción desaparece.

La corrección dejó, además, la mejor ilustración de la tesis anterior: al eliminar las
justificaciones equivocadas, \textbf{la tasa de anclaje bajó de 1,00 a 0,61 mientras la calidad
subía}, porque los falsos positivos pasaron de un texto fluido, anclado y falso a una descripción
factual generada por plantilla. Una métrica que empeora cuando el sistema mejora no puede
reportarse sola.

El trabajo deja, con la misma honestidad con que reporta sus logros, tres límites identificados y no
resueltos. El primero afecta a la propia señal que hizo posible el resultado principal: la
distinción entre un administrador legítimo y un atacante sobre el mismo servicio se apoya hoy en la
dirección de origen declarada en la configuración del cliente, que es suplantable, y cubre cuatro de
las seis categorías del caso de uso; las dos restantes exigen señales que el conjunto de datos no
contiene. El segundo es que el clasificador con ajuste fino, que resolvería ese límite de raíz,
permanece bloqueado porque el conjunto de datos disponible, concentrado en una sola familia de
ataque, no ofrece variedad suficiente para entrenarlo sin invalidar la propia evaluación. El tercero
es el alcance estadístico de la evaluación misma: dieciocho alertas soportadas, todas del mismo
activo y servicio, sostienen un resultado direccionalmente claro pero no permiten generalizar sus
cifras exactas a un cliente con un perfil de amenazas distinto.

\subsection*{Recomendaciones}

\begin{itemize}
  \item \textbf{Ampliar el conjunto de datos etiquetado con más familias de ataque y más activos},
        antes de intentar de nuevo el ajuste fino de un clasificador. Es la condición que hoy
        bloquea la pieza más importante que falta en el prototipo, y no depende de ninguna decisión
        externa a este proyecto.
  \item \textbf{Recuperar conocimiento adecuado a la clase, y no solo a la alerta.} Condicionar el
        enunciado a la clase ya bastó para eliminar las contradicciones, pero la recuperación sigue
        trayendo fichas sobre técnicas de ataque también cuando la alerta se ha descartado, y el
        corpus no contiene material que explique por qué algo \emph{no} es una amenaza. Ampliarlo en
        esa dirección es lo que permitiría que los falsos positivos recibieran una justificación
        generada en lugar de la plantilla.
  \item \textbf{Sustituir el modelo de \emph{embeddings} por uno dedicado} antes de intentar de
        nuevo ampliar el corpus. Está medido que ampliarlo con el modelo generativo como
        \emph{embedder} degrada la recuperación, y también que la consulta formulada por el modelo
        recupera peor que la consulta fija construida con los campos de la alerta. Las dos
        observaciones apuntan al mismo componente.
  \item \textbf{Resolver, junto con Domotai C.A., si el sistema de logs propietario que
        originalmente motivó el proyecto llegará a estar disponible}, y en qué formato entrega sus
        eventos. De confirmarse, se integraría como una segunda fuente a través del mismo módulo de
        ingesta que hoy consume las alertas de Wazuh, sin desplazarlo.
  \item \textbf{Completar la validación sobre un equipo de borde con firmware real}, hoy sustituido
        en el laboratorio por un contenedor Linux provisional: las acciones del catálogo están
        diseñadas para ese firmware, pero solo se han verificado sobre su sustituto.
  \item \textbf{Someter la calidad de la justificación a una segunda revisión manual independiente}
        antes de dar por definitiva la lectura de utilidad reportada en este informe, dado que el
        método de evaluación de esa dimensión no es completamente automatizable.
\end{itemize}

\subsection*{Trabajo futuro}

Quedan fuera del alcance de esta pasantía, y se documentan explícitamente para que no se
interpreten como un olvido: la respuesta automatizada sin validación humana previa, que el diseño
actual descarta de forma deliberada; la integración multicapa con fuentes de identidad, correo y
nube, que ampliaría la correlación más allá de la red del cliente; el aprendizaje continuo a partir
de la retroalimentación que el analista ya registra en cada validación, hoy almacenada pero no
utilizada para ajustar ningún modelo; y el despliegue y la escalabilidad del prototipo sobre
múltiples clientes de forma simultánea. Cada una de estas líneas es una extensión natural sobre la
arquitectura ya construida, y ninguna exige rediseñarla desde cero.
